\documentclass[12pt,letterpaper]{article}

\usepackage[utf8]{inputenc}
\usepackage[T1]{fontenc}
\usepackage[english]{babel}
\usepackage{float, enumitem, array, listings}
\usepackage{geometry, amsmath, amssymb, amsthm}
\usepackage{graphicx, xcolor, fancyhdr, fancyvrb}
\usepackage[normalem]{ulem}
\usepackage[hidelinks]{hyperref}

\geometry{margin=2cm, top=2cm, bottom=2cm}

\newcommand{\minic}{
  {M\footnotesize INI\normalsize C}
}

% ========================================
% CONSTANTES DEL ENCABEZADO
% ========================================
\newcommand{\encabezadoIzq}{Compiladores 2027-1}
\newcommand{\encabezadoCen}{Práctica 5}
\newcommand{\encabezadoDer}{\textit{Tabla de Símbolos y Ámbitos}}

% ========================================
% CONFIGURACION DEL ENCABEZADO
% ========================================
\pagestyle{fancy}
\fancyhf{}

\fancyhead[L]{\encabezadoIzq}
\fancyhead[C]{\encabezadoCen}
\fancyhead[R]{\encabezadoDer}

\fancyfoot[C]{\thepage}

\renewcommand{\headrulewidth}{0.1pt}
\renewcommand{\footrulewidth}{0pt}

% ========================================
% CONFIGURACION Y COLORES LSTLISTINGS
% ========================================
\renewcommand{\lstlistingname}{Código}

\definecolor{links}{HTML}{0D1894}

\definecolor{backcolour}{rgb}{0.93, 0.93, 0.93}
\definecolor{codegreen}{rgb}{0, 0.6, 0}
\definecolor{codegray}{rgb}{0.5, 0.5, 0.5}
\definecolor{codepurple}{rgb}{0.58, 0, 0.82}
\definecolor{keywordcolor}{rgb}{0.5, 0.0, 0.5}
\definecolor{codeorg}{HTML}{DB9400}
\definecolor{linkblue}{HTML}{1A167A}
\definecolor{redcolor}{HTML}{7A1805}
\definecolor{greencolor}{HTML}{0E8F10}
\definecolor{greeblucolor}{HTML}{0E6657}
\definecolor{orgcolor}{HTML}{B86D00}
\definecolor{highlightyellow}{HTML}{FFF176}
\definecolor{boxgray}{HTML}{F2F2F2}
\definecolor{borderblue}{HTML}{1A167A}

% ========================================
% CONFIGURACION LENGUAJES
% ========================================
\lstdefinestyle{modernc}{
    language=C,
    backgroundcolor=\color{backcolour},   
    commentstyle=\color{codegreen},
    keywordstyle=\color{keywordcolor}\bfseries,
    numberstyle=\tiny\color{codegray},
    stringstyle=\color{codepurple},
    basicstyle=\ttfamily\footnotesize,
    breakatwhitespace=false,
    breaklines=true,
    captionpos=b,
    keepspaces=true,
    numbers=left,
    numbersep=5pt,
    showspaces=false,                
    showstringspaces=false,
    showtabs=false,                  
    tabsize=4,
    frame=single,
    rulecolor=\color{black!20},
    frameround=tttt,
    emph={int,uint32_t,char,double,float,unsigned,void,bool,struct,size_t,
          Room, Client, pthread_mutex_t, NULL},
    emphstyle=\color{codeorg}\bfseries
}

\setlength{\emergencystretch}{3em}

\begin{document}

\thispagestyle{empty}
\begin{titlepage}
  \thispagestyle{empty}
  %--------------- COMANDO LINEA----------------->
  \newcommand{\linea}{\rule{\linewidth}{0.5mm}}                 
  \center
  %<<<<<<<<<<<<<<<<<<<<<<<<<<<<<<<<<<<<<<<<<<<

  %Add logos
  \includegraphics[width=0.24\textwidth]{../../unam_logo.png} \hspace{1.5cm}
  \includegraphics[width=0.25\textwidth]{../../fc_logo.png} \\[1cm]
  \textbf{\Large UNIVERSIDAD NACIONAL AUTÓNOMA}\\[3mm]
  \textbf{\Large DE MÉXICO} \\[6mm]
  \textit{\Large Facultad de Ciencias}

  \vfill

  %---------------TÍTULO----------------->
  \linea
  \vfill \vspace{1cm}
  \textbf{\Large Compiladores}\\[1.1cm]
  \textbf{\Large Tabla de símbolos y ámbitos}\\[0.9cm]
  \linea \\
  \vfill
  \vspace{0.7cm}
  %Title of the Research
  \textbf{\Large  Práctica 5}\\[6mm]
  %<<<<<<<<<<<<<<<<<<<<<<<<<<<<<<<<<<<<<<<<<<<

  %Team
  \vfill
  \textit{\small Presenta:}\\
  \textbf{\large  \textbf{Ramírez Juárez María Fernanda} - 321204747}\\[3mm]
  \textbf{\large  \textbf{Rojo Peña Manuel Ianluck} - 118005762}\\[3mm]

  %Profesor y Grupo
  \textit{\small Profesor:}\\
  \textbf{Pablo Martínez Yessica Janeth}\\[3mm]
  \textit{\small Ayudante de laboratorio:}\\
  \textbf{Salazar González Pedro Yamil}\\[3mm]
  \small Grupo: 7010, 2027-1\\[3mm]
  \vfill

  %Fecha or Date
  \textit{\small Fecha de entrega:}\\
         {\large 23 de octubre, 2026}
         
\end{titlepage}


%%%%%%%%%%%%%%%%%%%%%%%%%%%%%%%%%%%%%%%%%%
\newpage
%%%%%%%%%%%%%%%%%%%%%%%%%%%%%%%%%%%%%%%%%%

% ================================================================
\section*{Objetivo}
% ================================================================

\begin{quote}
  En esta práctica agregamos al compilador de \minic{} la primera etapa del análisis semántico. Hasta la
  Práctica 4 el compilador respondía \emph{¿qué elementos forman el programa?} y \emph{¿está sintácticamente
  bien escrito?}; ahora también responde \emph{¿tiene sentido lo que está escrito?}. Un programa puede
  cumplir la gramática y aun así usar una variable que nunca se declaró, o declarar dos veces el mismo
  identificador en el mismo ámbito.\\

  El analizador semántico no vuelve a leer el archivo fuente ni a pedir tokens, solo recorre el AST que produjo
  el parser. En esta entrega no verificamos tipos de expresiones, compatibilidad de asignaciones ni condiciones
  booleanas, y tampoco generamos código. Concretamente, el objetivo es:

  \begin{enumerate}[label=\arabic*.]
  \item Diseñar una tabla de símbolos que represente cada variable declarada
  \item Implementar ámbitos anidados, con un ámbito global y uno nuevo por cada bloque
  \item Verificar que toda variable se declare antes de usarse
  \item Prohibir la redeclaración dentro del mismo ámbito y permitir el \emph{shadowing} entre ámbitos distintos
  \item Reportar múltiples errores semánticos con línea y columna, sin detenerse en el primero
  \item Distinguir los errores semánticos del programa fuente de los fallos internos (memoria)
  \item Definir un único propietario para cada recurso y liberar la tabla sin fugas
  \end{enumerate}
\end{quote}

\bigskip

% ================================================================
\section*{Evolución respecto de la Práctica 4}
% ================================================================

\subsection*{Módulos modificados o agregados}

\begin{itemize}
\item \texttt{symbol.h} y \texttt{symbol.c}: módulo nuevo con el tipo \texttt{Type}, la estructura
  \texttt{Symbol} y sus funciones de creación y destrucción.
  
\item \texttt{scope.h} y \texttt{scope.c}: módulo nuevo con la estructura \texttt{Scope}, una tabla hash
  dinámica por ámbito y las operaciones de búsqueda e inserción dentro de un ámbito.
  
\item \texttt{symbol\_table.h} y \texttt{symbol\_table.c}: módulo nuevo con la pila de ámbitos, la entrada y
  salida de ámbitos, la búsqueda local y visible, la inserción y la destrucción.
  
\item \texttt{semantic.h} y \texttt{semantic.c}: módulo nuevo con el analizador, el recorrido del AST y los
  diagnósticos.
  
\item \texttt{main.c}: después del parser ejecuta el análisis semántico y cambia la salida normal.
  
\item \texttt{Makefile}: incorpora los archivos nuevos y las pruebas de esta práctica.
  
\item \texttt{lexer}, \texttt{token}, \texttt{keywords}, \texttt{buffer\_manage}, \texttt{parser.c},
  \texttt{grammar.c}, \texttt{ast.c} y \texttt{ast.h} se conservan sin cambios.
\end{itemize}

\subsection*{Arquitectura actualizada}

\begin{verbatim}
main
 |
 v
Lexer ----> Parser ----> AST
                          |
                          v
                  Semantic Analyzer
                   |-- recorrido del arbol
                   |-- manejo de scopes
                   |-- insercion y busqueda de simbolos
                   '-- diagnosticos
                          |
                          v
                    Symbol Table
                          |
                          v
                 Stack de Scopes (global -> ... -> actual)
\end{verbatim}

Como siempre, cada módulo tiene una responsabilidad única. El lexer reconoce tokens, el parser construye el AST,
el analizador semántico decide si el árbol tiene sentido, y la tabla de símbolos solo almacena símbolos y
administra ámbitos. Pero la tabla no conoce las reglas del lenguaje, no imprime diagnósticos ni recorre el AST.

\subsection*{Cambio de salida en \texttt{main.c}}

A partir de esta práctica la salida normal es y será diferente:

\begin{itemize}
\item Programa válido: \texttt{stdout} contiene exactamente \texttt{Programa semanticamente correcto.} y el
  código de salida es \texttt{0}.
  
\item Errores léxicos, sintácticos o semánticos: \texttt{stdout} queda vacío, los diagnósticos van a
  \texttt{stderr} y el código de salida es \texttt{1}.
  
\item Uso incorrecto, error de entrada y salida o fallo interno de memoria: se informa por \texttt{stderr} y el
  código de salida es \texttt{2}.
\end{itemize}

Ya no se imprimen el AST ni el mensaje \texttt{Programa sintacticamente correcto.} en la ejecución normal.
Para no perder la herramienta de depuración, \texttt{ast\_print} se conserva y se activa con la bandera
\texttt{--ast}, que reproduce exactamente la salida de la Práctica 4 y no ejecuta el análisis semántico.

\begin{lstlisting}[style=modernc, caption={Flujo de \texttt{compile} en \texttt{main.c}, resumido.}]
    result = parser_parse_program(&parser, &root);

    switch (result) {
    case PARSE_OK:
        if (show_ast) { /* modo depuracion: comportamiento de la P4 */
            printf("AST:\n");
            ast_print(root);
            printf("Programa sintacticamente correcto.\n");
            exit_code = MINIC_EXIT_OK;
        } else {
            exit_code = analyze(root);
        }
        break;
    case PARSE_INVALID:
        exit_code = MINIC_EXIT_INVALID; /* no se corre el semantico */
        break;
    default:
        exit_code = MINIC_EXIT_INTERNAL;
        break;
    }

    ast_destroy(root);
    parser_destroy(&parser);
    lexer_destroy(&lexer);
\end{lstlisting}

\bigskip

El orden de liberación es siempre el mismo, el analizador se destruye dentro de \texttt{analyze} (y después se
liberan el AST), el parser y el lexer. Como el analizador no es propietario de los nodos del árbol, el orden
entre el analizador y el AST no puede producir accesos inválidos.

% ================================================================
\section*{Diseño de la tabla de símbolos}
% ================================================================

\subsection*{Representación de un símbolo}

Cada variable declarada se representa con la estructura \texttt{Symbol}:

\begin{lstlisting}[style=modernc, caption={Estructura \texttt{Symbol} de \texttt{symbol.h}.}]
typedef enum { TYPE_INT, TYPE_BOOL, TYPE_ERROR, TYPE_UNKNOWN } Type;

typedef struct Symbol {
    char *name;           //copia propia
    Type type;
    int scope_level;
    int line;
    int column;
    ASTNode *declaration; //NO es propiedad del simbolo
    struct Symbol *next;  //encadenamiento en el bucket
} Symbol;
\end{lstlisting}

\bigskip

El símbolo conserva el nombre, el tipo declarado, el nivel de ámbito, la posición de la declaración y una
referencia al nodo del AST donde se declaró. En esta práctica solo se registran \texttt{TYPE\_INT} y
\texttt{TYPE\_BOOL}; \texttt{TYPE\_ERROR} y \texttt{TYPE\_UNKNOWN} quedan reservados para la verificación de
tipos. El campo \texttt{next} se agregó a la estructura recomendada para poder encadenar los símbolos que
caen en el mismo bucket de la tabla hash.

\begin{itemize}
\item \texttt{symbol\_create} reserva el símbolo y copia el nombre. Si cualquiera de las dos reservas falla,
  libera lo que ya había reservado y devuelve \texttt{NULL}. El nombre es una copia propia, por lo que no
  depende de la memoria del lexer ni del parser.
  
\item \texttt{symbol\_destroy} libera el nombre y la estructura, acepta \texttt{NULL} y \textbf{nunca} libera
  el nodo de declaración, porque ese nodo pertenece al AST.
\end{itemize}

La posición del símbolo es la que ya fija el AST de la Práctica 4, la declaración usa la posición del
\emph{identificador} (no la de \texttt{int} o \texttt{bool}), de modo que \texttt{int x;} produce la posición
\texttt{[1:4]}. Esto coincide con los diagnósticos que pide el enunciado.

\subsection*{Estructura de un ámbito}

\begin{lstlisting}[style=modernc, caption={Estructura \texttt{Scope} de \texttt{scope.h}.}]
typedef struct Scope {
    Symbol **buckets;
    size_t bucket_count;
    size_t symbol_count;
    int level;
    struct Scope *parent;
} Scope;
\end{lstlisting}

\bigskip

Cada ámbito guarda únicamente los símbolos declarados directamente en él, su nivel (el global es \texttt{0} y
cada bloque suma uno respecto de su padre) y un apuntador a su ámbito exterior. Para almacenar los símbolos
elegimos una tabla hash con encadenamiento:

\begin{itemize}
\item La función de dispersión es \texttt{djb2} sobre el nombre del identificador.
  
\item Cada ámbito comienza con 16 buckets.
  
\item Antes de insertar se verifica el factor de carga: si \texttt{(symbol\_count + 1) * 4 > bucket\_count * 3},
  es decir, si superaría 0.75, el arreglo de buckets duplica su tamaño y todos los símbolos se redistribuyen
  (\emph{rehash}).
  
\item Por tanto no existe un límite fijo de símbolos por ámbito, y las colisiones se resuelven con la lista
  encadenada de cada bucket.
\end{itemize}

\begin{lstlisting}[style=modernc, caption={Crecimiento de la tabla hash en \texttt{scope.c}.}]
static int
scope_grow(Scope *scope)
{
    size_t new_count = scope->bucket_count * 2;
    size_t i;
    Symbol **nb = calloc(new_count, sizeof *nb);
    if (nb == NULL)
        return 0;
    for (i = 0; i < scope->bucket_count; i++) {
        Symbol *s = scope->buckets[i];
        while (s != NULL) {
            Symbol *next = s->next;
            size_t idx = hash_name(s->name) % new_count;
            s->next = nb[idx];
            nb[idx] = s;
            s = next;
        }
    }
    free(scope->buckets);
    scope->buckets = nb;
    scope->bucket_count = new_count;
    return 1;
}
\end{lstlisting}

\bigskip

El nuevo arreglo se reserva primero y el anterior solo se libera cuando la redistribución terminó, de modo que
si la reserva falla la tabla original queda intacta. En ese caso \texttt{scope\_insert} devuelve un fallo
interno y no un duplicado.

\subsection*{La tabla de símbolos y la pila de ámbitos}

\begin{lstlisting}[style=modernc, caption={Estructura \texttt{SymbolTable} de \texttt{symbol\_table.h}.}]
typedef struct {
    Scope *current_scope;
    size_t scope_count;
} SymbolTable;
\end{lstlisting}

\bigskip

La pila de ámbitos se representa con la cadena de apuntadores \texttt{parent}: \texttt{current\_scope} es la
cima y el ámbito global es la base (su \texttt{parent} es \texttt{NULL}). No se necesita un arreglo aparte, y
entrar o salir de un ámbito solo mueve un apuntador.

\subsection*{Códigos de resultado}

Las operaciones que pueden fallar distinguen entre éxito, duplicado y fallo interno:

\begin{lstlisting}[style=modernc, caption={Códigos de resultado de la tabla y del analizador.}]
typedef enum {
    ST_OK,              //operacion correcta
    ST_DUPLICATE,       //el nombre ya existe en el ambito actual
    ST_INTERNAL_ERROR,  //fallo de memoria o estado invalido
    ST_AT_GLOBAL        //se intento salir del ambito global
} SymbolTableResult;

typedef enum {
    SEM_OK,             //sin errores semanticos
    SEM_ERRORS,         //errores semanticos en el programa fuente
    SEM_INTERNAL_ERROR  //fallo interno; analisis detenido
} SemanticResult;
\end{lstlisting}

\bigskip

Esta separación es importante pues un duplicado es un error del programa fuente que se puede recuperar y
acumular, mientras que un fallo de memoria es un error interno que detiene el análisis. Con un único código de
error no habría forma de evitar que un \texttt{malloc} fallido se reportara como una redeclaración.

\subsection*{Operaciones de la tabla}

\begin{itemize}
\item \texttt{symbol\_table\_init}: deja la tabla en un estado seguro (sin ámbito y con contador en cero) antes
  de reservar el ámbito global. Si la reserva falla, la tabla queda en un estado que \texttt{destroy} puede
  procesar sin riesgo.
  
\item \texttt{symbol\_table\_enter\_scope}: crea un ámbito hijo del actual y solo lo convierte en el ámbito
  vigente si la reserva terminó bien. Si falla, \texttt{current\_scope} no cambia.
  
\item \texttt{symbol\_table\_exit\_scope}: restaura el ámbito padre y libera el ámbito abandonado. Rechaza la
  salida del ámbito global devolviendo \texttt{ST\_AT\_GLOBAL}.
  
\item \texttt{symbol\_table\_lookup\_current}: busca únicamente en el ámbito actual.
  
\item \texttt{symbol\_table\_lookup}: búsqueda visible, del ámbito actual hacia el global.
  
\item \texttt{symbol\_table\_insert}: inserta solo en el ámbito actual y devuelve \texttt{ST\_OK},
  \texttt{ST\_DUPLICATE} o \texttt{ST\_INTERNAL\_ERROR}.
  
\item \texttt{symbol\_table\_destroy}: libera todos los ámbitos que todavía pertenecen a la tabla, desde el
  actual hasta el global, y es segura tras una inicialización incompleta.
\end{itemize}

\subsection*{Búsqueda local y búsqueda visible}

Son dos operaciones distintas con dos propósitos distintos:

\begin{itemize}
\item La \textbf{búsqueda local} (\texttt{lookup\_current}) revisa solo el ámbito actual y no consulta los
  exteriores. La usamos al \emph{declarar}, para detectar redeclaraciones.
  
\item La \textbf{búsqueda visible} (\texttt{lookup}) revisa el ámbito actual y, si no encuentra el nombre,
  sigue con el padre hasta llegar al global. Devuelve \texttt{NULL} si el identificador no existe en ninguno.
  Se usa al \emph{usar} un identificador, para validar que esté declarado.
\end{itemize}

\begin{lstlisting}[style=modernc, caption={Búsqueda visible en \texttt{symbol\_table.c}.}]
Symbol *
symbol_table_lookup(SymbolTable *table, const char *name)
{
    Scope *s;
    if (table == NULL)
        return NULL;
    for (s = table->current_scope; s != NULL; s = s->parent) {
        Symbol *found = scope_lookup(s, name);
        if (found != NULL)
            return found;
    }
    return NULL;
}
\end{lstlisting}

\bigskip

Como la búsqueda avanza del ámbito más interno al más externo, el primer símbolo encontrado es siempre el de
la declaración más cercana; eso es lo que implementa el \emph{shadowing}.

% ================================================================
\section*{Propiedad y liberación de memoria}
% ================================================================

Definimos un único propietario para cada recurso:

\begin{itemize}
\item \textbf{El símbolo} es propietario de la copia de su nombre. No es propietario del nodo de declaración
  del AST: \texttt{symbol\_destroy} nunca lo libera.
  
\item \textbf{El ámbito} es propietario de sus símbolos y de su arreglo de buckets.

\item \textbf{La tabla} es propietaria de los ámbitos que forman la pila vigente.

\item \textbf{El analizador} es propietario de la tabla, pero \emph{no} del AST: el árbol sigue perteneciendo a
  \texttt{main.c}, que lo libera con \texttt{ast\_destroy}.

\item \textbf{Inserción:} la tabla adquiere el símbolo únicamente si \texttt{symbol\_table\_insert} devuelve
  \texttt{ST\_OK}. En cualquier otro caso, quien llamó conserva la responsabilidad de liberarlo.
\end{itemize}

\subsection*{Estrategia de liberación de ámbitos}

Elegimos la \textbf{liberación inmediata}, al salir de un ámbito se destruyen sus símbolos y su arreglo de
buckets. Esta estrategia es válida en esta práctica porque el AST no guarda referencias a los símbolos
(\texttt{Symbol*}); la asociación entre usos y declaraciones se resuelve durante el recorrido y no se
conserva después. Además, deja la tabla con tamaño proporcional al anidamiento actual y no al del programa
completo.

En la Práctica 6, como comentamos en clase, esta estrategia tendría que cambiar: los ámbitos retirados deberían
moverse a una lista de ámbitos conservados y liberarse hasta \texttt{symbol\_table\_destroy}, para no dejar
apuntadores colgantes. Dejamos esta decisión documentada en un comentario de \texttt{symbol\_table.c}.

\begin{lstlisting}[style=modernc, caption={Salida de un ámbito.}]
SymbolTableResult
symbol_table_exit_scope(SymbolTable *table)
{
    Scope *old;
    if (table == NULL || table->current_scope == NULL)
        return ST_INTERNAL_ERROR;
    if (table->current_scope->parent == NULL)
        return ST_AT_GLOBAL;
    old = table->current_scope;
    table->current_scope = old->parent;
    table->scope_count--;
    scope_destroy(old);
    return ST_OK;
}
\end{lstlisting}

% ================================================================
\section*{Recorrido del AST}
% ================================================================

\subsection*{Estructura del analizador}

\begin{lstlisting}[style=modernc, caption={Estructura \texttt{SemanticAnalyzer} de \texttt{semantic.h}.}]
typedef struct {
    SymbolTable symbol_table;
    int error_count;
    int had_internal_error;
} SemanticAnalyzer;

int semantic_analyzer_init(SemanticAnalyzer *analyzer);
SemanticResult semantic_analyze(SemanticAnalyzer *analyzer, ASTNode *program);
void semantic_analyzer_destroy(SemanticAnalyzer *analyzer);
\end{lstlisting}

\bigskip

El analizador contabiliza los errores semánticos en \texttt{error\_count} y marca los fallos internos con
\texttt{had\_internal\_error}, de modo que ambos se pueden distinguir. \texttt{semantic\_analyze} devuelve
\texttt{SEM\_OK} si no hubo errores, \texttt{SEM\_ERRORS} si el programa fuente tiene errores semánticos y
\texttt{SEM\_INTERNAL\_ERROR} si ocurrió un fallo interno.

\subsection*{Funciones de visita}

El recorrido es recursivo y está dividido en funciones pequeñas, una por categoría de nodo, en lugar de una
sola función extensa. Todas son \texttt{static} dentro de \texttt{semantic.c}:

\begin{itemize}
\item \texttt{visit\_program}: recorre las sentencias en el orden en que aparecen y se detiene si hubo un
  fallo interno. El ámbito global ya existe desde \texttt{semantic\_analyzer\_init}.
  
\item \texttt{visit\_statement}: despacha según el tipo de nodo hacia la función correspondiente.

\item \texttt{visit\_block}: entra a un ámbito, recorre las sentencias y sale.

\item \texttt{visit\_variable\_declaration}: registra la variable y luego visita el inicializador.

\item \texttt{visit\_assignment}: valida el destino y visita el valor.

\item \texttt{visit\_print\_statement}, \texttt{visit\_if\_statement} y \texttt{visit\_while\_statement}:
  visitan sus expresiones y sentencias hijas.

\item \texttt{visit\_expression}: busca los identificadores dentro de expresiones binarias y unarias; los
  literales no requieren búsqueda.
\end{itemize}

Las declaraciones no se insertan por adelantado: se procesan en el orden del programa, de modo que una
declaración solo es visible después de haber sido recorrida. Por eso \texttt{x = 1; int x;} produce un error.

\subsection*{Declaraciones}

\begin{lstlisting}[style=modernc, caption={\texttt{visit\_variable\_declaration} (resumida).}]
    Symbol *prev = symbol_table_lookup_current(&a->symbol_table, name);

    if (prev != NULL) {
        /* No se sobrescribe el simbolo original. */
        report_redeclaration(a, prev, node->line, node->column);
    } else {
        Symbol *sym = symbol_create(name, type_from_ast(...),
            a->symbol_table.current_scope->level,
            node->line, node->column, node);
        if (sym == NULL) { report_internal_error(a); return; }
        r = symbol_table_insert(&a->symbol_table, sym);
        if (r != ST_OK) {          /* la tabla no adquirio el simbolo */
            symbol_destroy(sym);
            if (r == ST_INTERNAL_ERROR) {
                report_internal_error(a);
                return;
            }
        }
    }
    /* El simbolo ya es visible: int x = x; es valido. */
    visit_expression(a, node->data.variable_declaration.initializer);
\end{lstlisting}

\bigskip

Las decisiones de diseño relevantes son:

\begin{itemize}
\item La redeclaración se detecta con la búsqueda \emph{local}, por lo que declarar un nombre que existe solo
  en un ámbito exterior es válido (shadowing).

\item Ante una redeclaración, el símbolo original no se reemplaza. Esto es importante para la verificación de
  tipos de la siguiente práctica, y lo comprobamos con una prueba propia, en
  \texttt{int x; bool x; int x;} ambos errores señalan la declaración original en \texttt{[1:4]}; si el
  símbolo se hubiera sobrescrito, el segundo error apuntaría a \texttt{[2:5]}.

\item El símbolo se inserta \emph{antes} de visitar el inicializador. Siguiendo la regla de esta práctica
  (``la variable es visible a partir del momento en que su declaración fue registrada''), \texttt{int x = x;}
  no produce un error de variable no declarada. No se evalúa si el valor está inicializado.

\item Aun cuando la declaración sea una redeclaración, el inicializador se sigue visitando para poder
  detectar otros errores en la misma línea.

\item Si un símbolo recién creado no puede insertarse, la función que lo creó lo libera, porque la tabla solo
  adquiere símbolos insertados con \texttt{ST\_OK}.
\end{itemize}

\subsection*{Asignaciones, identificadores y expresiones}

Una asignación busca su destino con la búsqueda visible y reporta un error si no existe, y después recorre
la expresión del valor. Como el nodo de asignación ya guarda la posición del identificador destino,
\texttt{x = 5;} se reporta en \texttt{[1:0]}. Cada nodo identificador que representa un uso se valida con
una búsqueda visible; los identificadores que forman parte de una declaración no se validan como usos.\\

Las expresiones se recorren solo para localizar identificadores. No se infiere ningún tipo ni se comprueba la
compatibilidad entre el destino y el valor, tal como pide el enunciado: esa verificación corresponde a la
siguiente práctica.

\subsection*{Bloques}

\begin{lstlisting}[style=modernc, caption={\texttt{visit\_block}.}]
static void
visit_block(SemanticAnalyzer *a, ASTNode *node)
{
    size_t i;

    if (symbol_table_enter_scope(&a->symbol_table) != ST_OK) {
        report_internal_error(a);
        return;                    /* no se entro: no se sale */
    }
    for (i = 0; i < node->data.block.statements.count; i++) {
        if (a->had_internal_error)
            break;
        visit_statement(a, node->data.block.statements.items[i]);
    }
    if (symbol_table_exit_scope(&a->symbol_table) != ST_OK)
        report_internal_error(a);
}
\end{lstlisting}

\bigskip

La salida del ámbito se realiza siempre que la entrada tuvo éxito, incluso si se detectaron errores
semánticos dentro del bloque. Si la entrada falla, no se llama a la salida, porque no hay ámbito que
abandonar. Si ocurre un fallo interno dentro del bloque, el recorrido se detiene, el ámbito se abandona y el
fallo se propaga.

\subsection*{Condicionales y ciclos}

Para \texttt{if} se recorren la condición, la rama verdadera y la rama falsa cuando existe; para
\texttt{while}, la condición y el cuerpo. No se comprueba que la condición sea booleana, solo los nodos
\texttt{Block} crean ámbitos: una rama que es una sola sentencia sin llaves, como en
\texttt{if (c) x = 1; else x = 2;}, se recorre en el ámbito vigente.

% ================================================================
\section*{Diagnósticos y recuperación}
% ================================================================

\subsection*{Formato de los errores}

Los errores semánticos se envían a \texttt{stderr} con la categoría, la línea, la columna y una descripción que
incluye el nombre del identificador. La columna usa la convención del lexer (desde 0), igual que los errores
léxicos y sintácticos.

\begin{itemize}
\item \textbf{Variable no declarada:}
\begin{verbatim}
Error semantico [1:0]:
la variable 'x' no ha sido declarada.
\end{verbatim}
\item \textbf{Redeclaración,} con referencia a la declaración anterior:
\begin{verbatim}
Error semantico [2:5]:
la variable 'x' ya fue declarada
en este ambito.
Declaracion anterior [1:4].
\end{verbatim}
\end{itemize}

\subsection*{Acumulación de errores}

A diferencia de algunos errores sintácticos, un error semántico no impide seguir recorriendo el árbol. El
analizador incrementa \texttt{error\_count}, reporta el diagnóstico y continúa, por lo que detecta varios
errores en un mismo programa. Por ejemplo, para

\begin{verbatim}
x = 1;
y = 2;
int z;
int z;
print(w);
\end{verbatim}

se reportan cuatro errores: \texttt{x} y \texttt{y} no declaradas, la redeclaración de \texttt{z} y
\texttt{w} no declarada. El orden de los diagnósticos es el del recorrido, que coincide con el del programa.

\subsection*{Errores internos}

Los fallos de memoria al reservar símbolos, cadenas o ámbitos son errores internos, no errores semánticos.
\texttt{report\_internal\_error} imprime un diagnóstico por \texttt{stderr} y activa
\texttt{had\_internal\_error}; los recorridos de programa y bloque revisan esa bandera antes de cada
sentencia y se detienen de forma controlada. El resultado es \texttt{SEM\_INTERNAL\_ERROR}, el código de
salida es \texttt{2} y los recursos adquiridos se liberan en \texttt{semantic\_analyzer\_destroy}. Nunca se
informa una falla de memoria como redeclaración, porque \texttt{ST\_DUPLICATE} y \texttt{ST\_INTERNAL\_ERROR}
son resultados distintos.

% ================================================================
\section*{Resultados y pruebas}\label{sec:resultados}
% ================================================================

\subsection*{Cambios en el \texttt{Makefile}}

Respecto de la Práctica 4, el \texttt{Makefile} cambió en los siguientes puntos:

\begin{itemize}

\item Se agregaron los archivos \texttt{src/semantic/symbol.c}, \texttt{scope.c}, \texttt{symbol\_table.c} y
\texttt{semantic.c} a la variable \texttt{SOURCES}, y \texttt{include/ast/*.h} e
\texttt{include/semantic/*.h} a \texttt{HEADERS}, de modo que \texttt{make} recompile el ejecutable cuando
cambia cualquiera de ellos.

\item Se agregaron las variables que indican la ubicación de las pruebas públicas y propias de esta práctica:

\begin{verbatim}
PUBLIC_P5_INPUT_DIR = tests/public/p5/inputs
PUBLIC_P5_EXPECTED_DIR = tests/public/p5/expected

UNIT_SEMANTIC_INPUT_DIR = tests/unit_tests/semantic/inputs
UNIT_SEMANTIC_EXPECTED_DIR = tests/unit_tests/semantic/expected
\end{verbatim}

\item Se agregó el objetivo \texttt{test-unit-semantic}, que usa la plantilla \texttt{run\_test\_dir} y se
incluyó en \texttt{test-unit}. La plantilla compara siempre \texttt{stdout} con el archivo \texttt{.out} y,
cuando existen, también \texttt{stderr} con el \texttt{.err} y el código de salida con el \texttt{.code}.

\begin{verbatim}
test-unit-semantic: $(TARGET)
$(call run_test_dir,$(UNIT_SEMANTIC_INPUT_DIR),
  $(UNIT_SEMANTIC_EXPECTED_DIR),,unit-semantic)
\end{verbatim}

\item Se agregó el objetivo \texttt{test-p5} para las pruebas públicas de la práctica, incluido en
\texttt{test-cli}.

\item Como la salida normal cambió, \texttt{test-unit-parser} y \texttt{test-unit-ast} ahora invocan el
ejecutable con la bandera \texttt{--ast}, y los scripts de las pruebas públicas de las Prácticas 3 y 4
también. De esta forma conservan sus archivos esperados sin modificarlos.

\begin{verbatim}
test-unit-ast: $(TARGET)
$(call run_test_dir,$(UNIT_AST_INPUT_DIR),
  $(UNIT_AST_EXPECTED_DIR),--ast,unit-ast)
\end{verbatim}

\end{itemize}

\subsection*{Pruebas públicas de la Práctica 5}

Para ejecutar estas pruebas para la Práctica, es necesario utilizar el siguiente comando después de compilar:
\begin{verbatim}
sh tests/public/p5/run_public_tests.sh ./minic
\end{verbatim}

Al ejecutar las pruebas públicas proporcionadas, las 24 comprobaciones automáticas fueron aprobadas.

\begin{verbatim}
sh tests/public/p5/run_public_tests.sh ./minic
[OK automatico] p01_empty
[OK automatico] p02_literals
[OK automatico] p03_declarations
[OK automatico] p04_valid_use
[OK automatico] p05_undeclared_assignment
[OK automatico] p06_use_before_declaration
[OK automatico] p07_global_redeclaration
[OK automatico] p08_local_redeclaration
[OK automatico] p09_nested_visibility
[OK automatico] p10_out_of_scope
[OK automatico] p11_shadowing
[OK automatico] p12_sibling_scopes
[OK automatico] p13_sibling_visibility
[OK automatico] p14_initializer_use
[OK automatico] p15_self_initializer
[OK automatico] p16_multiple_errors
[OK automatico] p17_condition_uses
[OK automatico] p18_if_else
[OK automatico] p19_while
[OK automatico] p20_types_deferred
[OK automatico] p21_syntax_error
[OK automatico] p22_lexical_error
[OK automatico] p23_expression_uses
[OK automatico] p24_empty_blocks

Comprobaciones automaticas: 24 aprobadas, 0 fallidas.
Casos que requieren revisar diagnosticos: 12.
\end{verbatim}

El script marca doce casos como \emph{REVISAR}, pues sus diagnósticos deben compararse contra
\texttt{expected/CASOS.md}. Revisamos algunos de manera manual y los mensajes coinciden, por ejemplo, en
\texttt{p07\_global\_redeclaration} el diagnóstico es:

\begin{verbatim}
Error semantico [2:5]:
la variable 'x' ya fue declarada
en este ambito.
Declaracion anterior [1:4].
\end{verbatim}

y en \texttt{p16\_multiple\_errors} se reportan, en orden, los cuatro errores esperados (\texttt{[1:0]},
\texttt{[2:0]}, \texttt{[4:4]} con declaración anterior en \texttt{[3:4]}, y \texttt{[5:6]}). Los casos
\texttt{p21\_syntax\_error} y \texttt{p22\_lexical\_error} comprueban que, ante un error sintáctico o léxico,
el análisis semántico no se ejecuta y se conserva el diagnóstico de la etapa correspondiente.

\subsection*{Pruebas propias del analizador semántico}

Además de las pruebas públicas, diseñamos 40 pruebas propias, con el mismo esquema de las pruebas del AST,
cada programa \texttt{.mc} tiene su salida estándar esperada (\texttt{.out}), sus diagnósticos esperados
(\texttt{.err}) y su código de salida (\texttt{.code}). Se organizaron en \texttt{tests/unit\_tests/semantic/}
con los prefijos \texttt{ut\_sem\_}.

Para no calcular los archivos esperados a mano, escribimos un modelo de referencia independiente de la
implementación en C (un analizador léxico y un recorrido con pila de ámbitos) y generamos los \texttt{.err}
con él. Antes de usarlo lo validamos contra seis casos de los ejemplos públicos
(\texttt{p07}, \texttt{p08}, \texttt{p14}, \texttt{p16}, \texttt{p17} y \texttt{p23}), con resultados idénticos.
De este modo las pruebas comparan la implementación en C contra un modelo distinto.

\begin{table}[H]
\centering
\small
\begin{tabular}{|p{3.4cm}|p{12cm}|}
\hline
\textbf{Tema} & \textbf{Casos} \\ \hline
Básicos & \texttt{empty}, \texttt{no\_declarations}, \texttt{one\_global}, \texttt{multiple\_globals} \\ \hline
Declaración y uso & \texttt{use\_after\_declaration}, \texttt{use\_before\_declaration},
  \texttt{undeclared\_completely}, \texttt{assignment\_target}, \texttt{expression\_use},
  \texttt{print\_use}, \texttt{condition\_use}, \texttt{statement\_order} \\ \hline
Redeclaración & \texttt{global\_redeclaration}, \texttt{block\_redeclaration},
  \texttt{redeclaration\_other\_type}, \texttt{redeclaration\_keeps\_original},
  \texttt{redeclaration\_initializer}, \texttt{long\_identifier\_redecl} \\ \hline
Shadowing y visibilidad & \texttt{shadow\_global}, \texttt{shadow\_multilevel}, \texttt{outer\_access},
  \texttt{global\_from\_nested}, \texttt{initializer\_uses}, \texttt{self\_initializer} \\ \hline
Salida de ámbito & \texttt{out\_of\_scope}, \texttt{out\_of\_scope\_nested},
  \texttt{sibling\_visibility}, \texttt{while\_local\_escape} \\ \hline
Bloques y control & \texttt{empty\_blocks}, \texttt{sibling\_blocks}, \texttt{if\_scopes},
  \texttt{if\_else\_valid}, \texttt{while\_locals}, \texttt{if\_without\_blocks\_ok},
  \texttt{deep\_nesting} \\ \hline
Recuperación & \texttt{multiple\_errors}, \texttt{errors\_inside\_blocks} \\ \hline
Estrés & \texttt{many\_symbols}, \texttt{many\_symbols\_redecl}, \texttt{long\_identifiers} \\ \hline
\end{tabular}
\caption{Casos propios de \texttt{tests/unit\_tests/semantic/inputs}, todos con prefijo \texttt{ut\_sem\_}.}
\end{table}

Algunos casos merecen explicación:

\begin{itemize}
\item \texttt{redeclaration\_keeps\_original} comprueba que una redeclaración no reemplaza al símbolo
  original: con \texttt{int x; bool x; int x;}, ambos errores deben apuntar a la declaración anterior
  \texttt{[1:4]}.

\item \texttt{many\_symbols} declara 300 variables globales y vuelve a declarar la mitad en un bloque, lo que
  obliga a la tabla hash a hacer varios \emph{rehash} y a resolver colisiones. \texttt{many\_symbols\_redecl}
  agrega una redeclaración de \texttt{v150} y verifica que la tabla la detecte entre cientos de símbolos.

\item \texttt{long\_identifiers} usa identificadores de 100 caracteres, para comprobar que el nombre se copia
  completo.

\item \texttt{statement\_order} verifica que las sentencias se procesan en orden: usar una variable antes de
  su declaración es un error, aunque la declaración aparezca después.

\item \texttt{errors\_inside\_blocks} y \texttt{multiple\_errors} verifican que el análisis continúa tras el
  primer error y que el ámbito se abandona correctamente aunque haya errores dentro.
\end{itemize}

\subsubsection*{Ejecución de las pruebas}

Para ejecutar nuestras pruebas de esta Práctica, es necesario utilizar el siguiente comando después de compilar:
\begin{verbatim}
$ make test-unit-semantic
\end{verbatim}

El sistema ejecuta cada archivo \texttt{.mc} y compara la salida estándar, los diagnósticos y el código de
salida contra los archivos esperados. Una prueba se considera correcta solo si los tres coinciden.

\begin{verbatim}
==== unit-semantic ====
PASS: ut_sem_assignment_target
PASS: ut_sem_block_redeclaration
PASS: ut_sem_condition_use
PASS: ut_sem_deep_nesting
PASS: ut_sem_empty_blocks
PASS: ut_sem_empty
PASS: ut_sem_errors_inside_blocks
PASS: ut_sem_expression_use
PASS: ut_sem_global_from_nested
PASS: ut_sem_global_redeclaration
PASS: ut_sem_if_else_valid
PASS: ut_sem_if_scopes
PASS: ut_sem_if_without_blocks_ok
PASS: ut_sem_initializer_uses
PASS: ut_sem_long_identifier_redecl
PASS: ut_sem_long_identifiers
PASS: ut_sem_many_symbols
PASS: ut_sem_many_symbols_redecl
PASS: ut_sem_multiple_errors
PASS: ut_sem_multiple_globals
PASS: ut_sem_no_declarations
PASS: ut_sem_one_global
PASS: ut_sem_outer_access
PASS: ut_sem_out_of_scope
PASS: ut_sem_out_of_scope_nested
PASS: ut_sem_print_use
PASS: ut_sem_redeclaration_initializer
PASS: ut_sem_redeclaration_keeps_original
PASS: ut_sem_redeclaration_other_type
PASS: ut_sem_self_initializer
PASS: ut_sem_shadow_global
PASS: ut_sem_shadow_multilevel
PASS: ut_sem_sibling_blocks
PASS: ut_sem_sibling_visibility
PASS: ut_sem_statement_order
PASS: ut_sem_undeclared_completely
PASS: ut_sem_use_after_declaration
PASS: ut_sem_use_before_declaration
PASS: ut_sem_while_local_escape
PASS: ut_sem_while_locals
----
Total: 40  Pass: 40  Fail: 0]
\end{verbatim}

% ================================================================
\section*{Problemas encontrados}\label{sec:problemas}
% ================================================================

\subsection*{Las pruebas de la Práctica 4 dejaron de pasar}

Al cambiar la salida normal, la mayoría de las pruebas de la Práctica 4 fallaron, esperaban \texttt{AST:}, el
árbol y \texttt{Programa sintacticamente correcto.}, pero el programa ahora imprime
\texttt{Programa semanticamente correcto.}. Solo seguían pasando los casos con errores léxicos o sintácticos,
cuya salida no cambió. No era un error del compilador sino una consecuencia esperada del enunciado.
Lo resolvimos conservando la salida anterior bajo la bandera \texttt{--ast} y haciendo que los scripts y los
objetivos del \texttt{Makefile} de las prácticas anteriores la usen.

\subsection*{Duplicado frente a fallo interno}

Una primera idea es que la inserción devuelva solo ``éxito'' o ``fallo''. Con ese diseño, un \texttt{malloc}
fallido durante el \emph{rehash} se habría interpretado como una redeclaración. Por eso usamos tres
resultados distintos (\texttt{ST\_OK}, \texttt{ST\_DUPLICATE}, \texttt{ST\_INTERNAL\_ERROR}), y establecimos
que la tabla solo adquiere el símbolo cuando devuelve \texttt{ST\_OK}.

\subsection*{Propiedad del símbolo en la inserción}

Como la tabla no adquiere el símbolo cuando la inserción falla, \texttt{visit\_variable\_declaration} debe
liberarlo en esa ruta. Olvidarlo habría producido una fuga justo en el caso de un error poco frecuente, que
las pruebas normales no recorren.

\subsection*{Posición de la declaración}

Los diagnósticos del enunciado esperan que una declaración se reporte en la posición del identificador y no
en la de \texttt{int}. Revisamos \texttt{parse\_declaration} y confirmamos que, desde la Práctica 4, el nodo
ya almacenaba la posición del identificador, por lo que no fue necesario modificar el parser.

\subsection*{Límite por anidamiento profundo}

Como en las prácticas anteriores, el recorrido recursivo consume pila en proporción al anidamiento. Una
entrada con una cantidad extrema de bloques anidados puede agotarla. No agregamos un límite de profundidad y
lo dejamos registrado como limitación conocida.

% ================================================================
\section*{Conclusiones}
% ================================================================
Esta práctica nos mostró que el análisis semántico es, en gran parte, una cuestión de recordar lo que se ha
visto mientras se recorre el árbol. El AST de la práctica anterior ya tenía toda la información necesaria
(nombres, posiciones y estructura de bloques), así que no tuvimos que modificar el lexer, el parser ni el
árbol. Con essto seguimos confirmando que separar el lexer, el parser y el AST nos permite agregar una etapa
nueva sin tocar las anteriores.\\

La pila de ámbitos resuelve de una sola vez la visibilidad, la salida de ámbito y el \emph{shadowing}. La
búsqueda avanza del ámbito más interno al global, por lo que siempre se encuentra la declaración más
cercana, y al abandonar un bloque sus símbolos desaparecen sin ningún tratamiento adicional. La diferencia
entre la búsqueda local, para declarar, y la búsqueda visible, para usar, es lo que permite al mismo tiempo
prohibir la redeclaración en el mismo ámbito y permitirla en ámbitos distintos.\\

La parte más delicada fue, de nuevo, la propiedad de la memoria. Definir que la tabla adquiere el símbolo solo
si la inserción termina con \texttt{ST\_OK}, que el símbolo no es propietario del nodo del AST y que el
analizador no es propietario del árbol nos obligó a razonar sobre cada ruta de error. También fue útil separar
los errores semánticos, que se acumulan, de los fallos internos, que detienen el análisis de forma controlada.
Con esa distinción, el compilador nunca presenta un fallo de memoria como un error del programa fuente.\\

Las pruebas propias nos dieron una verificación independiente, al generar los resultados esperados con un modelo
de referencia distinto de la implementación en C, evitamos que un mismo error de comprensión estuviera en ambos
lados. La tabla de símbolos conserva ya el tipo declarado de cada variable y la declaración original después de
una redeclaración, por lo que la siguiente práctica podrá apoyarse en ella para inferir y verificar tipos.
Como limitación, el diseño sigue siendo recursivo y la liberación inmediata de ámbitos tendrá que revisarse
si el AST se anota con apuntadores a símbolos.

% ================================================================
\section*{División del trabajo}
% ================================================================
Para esta práctica, el trabajo tanto del código como del reporte fue realizado en conjunto por ambos miembros
del equipo durante tiempos libres dentro de la facultad. Discutiendo y analizando lo necesario a implementar,
diseñar y cómo realizarlo.

% ================================================================
\section*{Uso de herramientas de Inteligencia Artificial}
% ================================================================
Para la realización de esta práctica se utilizó la herramienta de Inteligencia Artificial ClaudeAI como apoyo
en la redacción del \texttt{CHANGELOG} debido al tiempo limitado que se contaba para realizar. Su uso se enfocó
principalmente en utilziar este reporte y seguir el formato del \texttt{CHANGELOG} echo en prácticas anteriores.
Se utilizó el siguiente prompt:

\begin{verbatim}
Utilizando este formato para el CHANGELOG de mis prácticas, y con este reporte de la práctica 5, dame el
fragmento a incluir para el CHANGELOG de la respectiva práctica 5. Mantén la información técnica y la escencia
original, pero haz que el texto sea más breve y claro sin perder el punto del reporte y cambios. No agregues
 información que no esté presente en el texto original.
\end{verbatim}

Por brevedad no incluímos el texto copiado, pero simplemente fue el contendio previo del \texttt{CHANGELOG} y el
texto del reporte.tex.
\begin{thebibliography}{99}
\bibitem{ref1}
  cppreference.com. (n.d.). C language. Disponible en
  \textcolor{links}{\uline{\url{https://en.cppreference.com/c/language}}}

\bibitem{ref2}
  cppreference.com. (n.d.). C11. Disponible en
  \textcolor{links}{\uline{\url{https://en.cppreference.com/c/11}}}

\bibitem{ref3}
  Free Software Foundation. (n.d.). GNU make. Disponible en
  \textcolor{links}{\uline{\url{https://www.gnu.org/software/make/manual/make.html}}}
\end{thebibliography}

\end{document}
